\chapter{Vectors and Lists}
\label{chap:vectors-lists}

\chapterguide
  {You should be able to create variables, call functions, and use logical operators.}
  {construct, index, sort, and calculate with vectors; and create, name, access, and modify lists.}

\section{R vectors}

A vector stores values of a common type and can be created with \texttt{:}, \texttt{seq()}, or \texttt{c()}.

\subsection{Constructing vectors}

\begin{lstlisting}[style=dsr]
# Integer-like sequence
v <- 5:13
print(v)

# Sequence beginning at a decimal value
v <- 6.6:12.6
print(v)

# End value does not fall exactly on the sequence
v <- 3.8:11.4
print(v)

# Explicit step size
print(seq(5, 9, by = 0.4))
\end{lstlisting}

Mixing character, numeric, and logical values triggers type coercion:

\begin{lstlisting}[style=dsr]
s <- c('apple', 'red', 5, TRUE)
print(s)
\end{lstlisting}

\begin{keyidea}
All elements in the mixed vector become character values. Use a list when components must retain different types.
\end{keyidea}

\subsection{Accessing vector elements}

The days vector demonstrates positive, logical, and negative indexing:

\begin{lstlisting}[style=dsr]
t <- c("Sun", "Mon", "Tue", "Wed", "Thurs", "Fri", "Sat")

# Position indexing
u <- t[c(2, 3, 6)]
print(u)

# Logical indexing
v <- t[c(TRUE, FALSE, FALSE, FALSE, FALSE, TRUE, FALSE)]
print(v)

# Negative indexing
x <- t[c(-2, -5)]
print(x)

# Select the seventh element
y <- t[7]
print(y)
\end{lstlisting}

Positive indices select positions, logical indices select TRUE positions, and negative indices exclude positions.

\begin{pitfall}
Numeric 0/1 indexing is not a Boolean mask in R: zeros are ignored and each 1 selects position 1. Use TRUE/FALSE values or explicit positions instead.
\end{pitfall}

\section{Vector arithmetic}

Arithmetic on equal-length vectors is element-wise:

\begin{lstlisting}[style=dsr]
v1 <- c(3, 8, 4, 5, 0, 11)
v2 <- c(4, 11, 0, 8, 1, 2)

print(v1 + v2)
print(v1 - v2)
print(v1 * v2)
print(v1 / v2)
\end{lstlisting}

Division by zero produces a non-finite result such as \texttt{Inf} or \texttt{NaN}.

\subsection{Vector element recycling}

Lab 3a next shortens the second vector:

\begin{lstlisting}[style=dsr]
v1 <- c(3, 8, 4, 5, 0, 11)
v2 <- c(4, 11)

# Handout note: v2 behaves as c(4, 11, 4, 11, 4, 11)
print(v1 + v2)
print(v1 - v2)
\end{lstlisting}

R recycles the shorter vector. A warning occurs when the longer length is not a multiple of the shorter length.

\section{Sorting vectors}

\begin{lstlisting}[style=dsr]
v <- c(3, 8, 4, 5, 0, 11, -9, 304)
print(sort(v))
print(sort(v, decreasing = TRUE))

v <- c("Red", "Blue", "yellow", "violet")
print(sort(v))
print(sort(v, decreasing = TRUE))
\end{lstlisting}

\texttt{sort()} handles numeric and character vectors; \texttt{decreasing = TRUE} reverses the order.

\section{R lists}

A list is introduced using a mixture of strings, numbers, a vector, and a logical value:

\begin{lstlisting}[style=dsr]
list_data <- list(
  "Red",
  "Green",
  c(21, 32, 11),
  TRUE,
  51.23,
  119.1
)
print(list_data)
\end{lstlisting}

Unlike a vector, a list preserves heterogeneous component types.

\subsection{Naming list elements}

\begin{lstlisting}[style=dsr]
list_data <- list(
  c("Jan", "Feb", "Mar"),
  list("green", 12.3)
)

names(list_data) <- c("1st Quarter", "A Inner list")
print(list_data)
\end{lstlisting}

Names make list access more readable.

\subsection{Accessing list elements}

Lists support position- and name-based access:

\begin{lstlisting}[style=dsr]
list_data <- list(
  c("Jan", "Feb", "Mar"),
  list("green", 12.3)
)
names(list_data) <- c("1st_Quarter", "A_Inner_list")

print(list_data[1])
print(list_data[2])
print(list_data$A_Inner_list)
print(which(list_data$`1st_Quarter` == "Feb"))
\end{lstlisting}

\texttt{which()} returns the position of \texttt{"Feb"} within the named component.

\begin{keyidea}
Single brackets return a sublist; double brackets or \texttt{\$} return the component itself.
\end{keyidea}

\subsection{Manipulating a list}

\begin{lstlisting}[style=dsr]
list_data <- list(
  c("Jan", "Feb", "Mar"),
  list("green", 12.3)
)
names(list_data) <- c("1st_Quarter", "A_Inner_list")

# Add a third element
list_data[3] <- "New_element"
print(list_data[3])

# Remove it
list_data[3] <- NULL
print(list_data[3])

# Update the second element
list_data[2] <- "updated_element"
print(list_data[2])
\end{lstlisting}

Assignment adds or updates an element; assigning \texttt{NULL} removes it.

\begin{quickcheck}
\begin{enumerate}
  \item What happens to the numeric and logical values in the handout's mixed character vector?
  \item How does negative vector indexing differ from positive position indexing?
  \item What repeated form of \texttt{v2} is stated in the vector-recycling example?
  \item Why is a list suitable for the Lab 3a example containing strings, numbers, a vector, and a logical value?
  \item How does \texttt{list\_data\$A\_Inner\_list} differ in intent from accessing an unnamed position?
\end{enumerate}
\end{quickcheck}

\begin{chapterrecap}
\begin{itemize}
  \item Vectors use a common type and support positive, logical, and negative indexing.
  \item Vector arithmetic is element-wise and may recycle shorter operands.
  \item Lists preserve mixed types and support named access and modification.
\end{itemize}
\end{chapterrecap}
